\documentclass[11pt]{article}
\usepackage{mathblatt}
% gesetzt von werkzeuge/setzer.py; Lösungen nach bau/bauregeln.md „Lösungen“.
\newcommand{\kennung}{4NT}
\begin{document}
\pfheftstil
\fancyfoot[C]{{\footnotesize\color{mbgrau}\kennung\hfill\thepage}}
\pfheftkopf{Prozentsatz berechnen}{Klasse 7 \textperiodcentered{} Lösungen}

\begin{pfloesung}{}
\lz{1b)}{$40\,\%$}{$\frac{20}{50} = \frac{40}{100}$}{}
\lz{c)}{$70\,\%$}{$\frac{140}{200} = \frac{70}{100}$}{}
\end{pfloesung}
\begin{pfloesung}{}
\lz{2b)}{$\frac{90}{600}$}{Ganzes: alle Lose}{}
\lz{c)}{$\frac{3}{16}$}{Ganzes $13 + 3 = 16$, nicht $13$}{}
\lz{d)}{$\frac{20}{80}$}{Rabatt $80 - 60 = 20$; Ganzes: alter Preis}{}
\lz{e)}{$\frac{45}{150}$}{Ganzes $45 + 105 = 150$}{}
\lz{f)}{$\frac{6}{24}$}{frei $24 - 18 = 6$}{}
\end{pfloesung}
\begin{pfloesung}{}
\lz{3b)}{$38\,\%$}{$\frac{19}{50} = \frac{38}{100}$ (mal $2$)}{}
\lz{c)}{$36\,\%$}{$\frac{9}{25} = \frac{36}{100}$ (mal $4$)}{}
\lz{d)}{$30\,\%$}{$\frac{3}{10} = \frac{30}{100}$ (mal $10$)}{}
\lz{e)}{$42\,\%$}{$\frac{84}{200} = \frac{42}{100}$ (durch $2$)}{}
\lz{f)}{$6\,\%$}{$\frac{18}{300} = \frac{6}{100}$ (durch $3$)}{}
\end{pfloesung}
\begin{pfloesung}{}
\lz{4b)}{$40\,\%$}{$18 : 45 = 0{,}4$}{}
\lz{c)}{$15\,\%$}{$51 : 340 = 0{,}15$}{}
\lz{d)}{$65\,\%$}{$117 : 180 = 0{,}65$}{}
\lz{e)}{$75\,\%$}{$9 : 12 = 0{,}75$}{}
\lz{f)}{unter 100\,\%}{Teil kleiner als Ganzes, also unter $100\,\%$; Ben rechnet Ganzes : Teil}{}
\end{pfloesung}
\begin{pfloesung}{}
\lz{5b)}{$54{,}2\,\%$}{$13 : 24 = 0{,}5416\ldots = 54{,}16\ldots\,\%$}{}
\lz{c)}{$74{,}5\,\%$}{$41 : 55 = 0{,}7454\ldots = 74{,}54\ldots\,\%$}{}
\lz{d)}{$36{,}7\,\%$}{$11 : 30 = 0{,}3666\ldots = 36{,}66\ldots\,\%$}{}
\lz{e)}{$77{,}8\,\%$}{$7 : 9 = 0{,}7777\ldots = 77{,}77\ldots\,\%$}{}
\end{pfloesung}
\begin{pfloesung}{}
\lz{6.}{a) $42{,}9\,\%$; b) $7{,}1\,\%$}{Ganzes $9 + 12 + 5 + 2 = 28$; $12 : 28 = 0{,}4285\ldots = 42{,}85\ldots\,\%$; $2 : 28 = 0{,}0714\ldots = 7{,}14\ldots\,\%$}{}
\end{pfloesung}
\begin{pfloesung}{}
\lz{7.}{$20\,\%$ Rabatt}{Rabatt $240 - 192 = 48\,€$; $48 : 240 = 0{,}2$ (nicht $192 : 240 = 80\,\%$)}{}
\end{pfloesung}
\begin{pfloesung}{}
\lz{8a)}{Nein, nur etwa $46{,}7\,\%$}{$14 : 30 = 0{,}4666\ldots = 46{,}66\ldots\,\%$; weniger als $50\,\%$}{}
\lz{b)}{Ja, etwa $41{,}7\,\%$}{Ganzes: abgestimmt haben $14 + 10 = 24$; $10 : 24 = 0{,}4166\ldots = 41{,}66\ldots\,\%$}{}
\end{pfloesung}
\begin{pfloesung}{}
\lz{9.}{$15{,}2\,\%$}{nicht verwandelt $46 - 39 = 7$; $7 : 46 = 0{,}1521\ldots = 15{,}21\ldots\,\%$}{}
\end{pfloesung}
\end{document}
